% ---------------------------------------------------------------------------
% MOTIVATING FIGURE (Sec. sec-intro). Self-contained view of the five-claim
% Voyager 1 pair that Sec. sec-example later works formally and App. F.1 gives
% in full. Four bands, top to bottom:
%   (a) the two responses, abridged to their claim- and relation-bearing spans;
%   (b) the relation graphs those spans induce (transcribed from
%       figures/graph-coherence-pair.tex, so the two figures cannot diverge);
%   (c) the per-claim posterior marginals, with each claim's prior ticked;
%   (d) the scores and the gap.
% Every number is produced by scripts/lcs_worked_pair.py from
% data/lcs/example-7-coherence-pair.json -- exact, 2^5 = 32-world enumeration,
% computed twice by independent paths and asserted equal. Do not edit a value
% here without rerunning that script.
% Node, edge and relw=p styles live in preamble.tex, shared with the Sec. 3.7
% and App. F figures, so no panel can drift from another. The two graph panels
% are drawn at ONE fixed scale (\mtvscale) rather than each in its own
% \resizebox, so a difference in bounding box cannot make one look larger.
% ---------------------------------------------------------------------------
\begin{figure}[p]
\centering
\footnotesize
\setlength{\tabcolsep}{0pt}
% one shared scale for both relation graphs
\newcommand{\mtvx}{13.2mm}
\newcommand{\mtvy}{11.6mm}

%% ===== band (0): the query both responses answer =========================
% Grey and unadorned: it belongs to neither response, and the point of showing
% it is that BOTH answer it. Its two parts are why a good answer must both
% assert a3 and dispatch the a5 confusion.
\colorbox{black!5}{\parbox{\dimexpr\textwidth-2\fboxsep\relax}{\centering
  \scriptsize \textbf{Query}\ \ \emph{Where is Voyager~1 now, and has it left
  the Solar System?}\strut}}

\smallskip
%% ===== band (a): the two responses ======================================
% Both boxes are set to ONE height (\mtvbox), so the pair aligns top and bottom
% however the two texts happen to break. Raise it if either box overflows.
\newcommand{\mtvbox}{33.5mm}
\begin{tabularx}{\textwidth}{@{}X@{\hspace{7pt}}X@{}}
\begin{minipage}[t]{\linewidth}
  \colorbox{blue!4}{\parbox[t][\mtvbox][t]{\dimexpr\linewidth-2\fboxsep\relax}{%
  \textbf{\textcolor{blue!55!black}{Response A}} \; \emph{coherent}\\[2pt]
  \scriptsize
  Voyager~1 crossed the heliopause in August 2012,\,\claimref{1} and in the days
  that followed its instruments measured a sharp rise in plasma
  density\,\claimref{2} --- \emph{the measurement that established the boundary
  had been passed.} Voyager~1 is therefore in interstellar space.\,\claimref{3}
  Two persistent misconceptions are worth clearing up. The first is that the
  plasma wave instrument had failed permanently before the
  crossing;\,\claimref{4} \emph{it cannot have, since it is precisely that
  instrument which returned the density profile the conclusion rests on.} The
  second is that Voyager~1 has left the Solar System
  altogether.\,\claimref{5} \emph{Being in interstellar space is not the same
  thing: the Sun's gravitational reach extends through the Oort cloud.}\strut}}
\end{minipage}
&
\begin{minipage}[t]{\linewidth}
  \colorbox{red!4}{\parbox[t][\mtvbox][t]{\dimexpr\linewidth-2\fboxsep\relax}{%
  \textbf{\textcolor{red!60!black}{Response B}} \; \emph{incoherent}\\[2pt]
  \scriptsize
  Voyager~1 crossed the heliopause in August 2012,\,\claimref{1} and after the
  crossing its instruments measured a sharp rise in plasma
  density.\,\claimref{2} \emph{That said, the probe cannot really be counted as
  being in interstellar space.}\,\claimref{3} The plasma wave instrument had
  \textbf{in fact} failed permanently before the crossing,\,\claimref{4}
  \emph{and it is precisely that failure which tells us the probe reached
  interstellar space.} \emph{The instrument failure also means} Voyager~1 has
  left the Solar System entirely.\,\claimref{5} \emph{Of course, the crossing
  rules out any such failure, and the density measurement is inconsistent with
  the probe having left the Solar System.}\strut}}
\end{minipage}
\end{tabularx}

\smallskip
{\scriptsize
  \textbf{The same five claims} $a_1,\dots,a_5$ (\claimref{i} marks each first
  mention) $\Rightarrow$ \textbf{identical per-claim factuality verdicts.}
  Only the italicized \emph{relation-bearing} spans differ.}

\medskip
%% ===== band (b): the relation graphs, at one shared scale ===============
\begin{tabularx}{\textwidth}{@{}X@{\hspace{7pt}}X@{}}
\begin{minipage}[b]{\linewidth}\centering
\begin{tikzpicture}[x=\mtvx, y=\mtvy]
  \node[atomtrue]  (a1) at (0,0)        {$a_1$};
  \node[atomtrue]  (a2) at (1.15,0)     {$a_2$};
  \node[atomtrue]  (a3) at (2.3,0)      {$a_3$};
  \node[atomfalse] (a4) at (1.7,-1.55)  {$a_4$};
  \node[atomfalse] (a5) at (2.9,-1.55)  {$a_5$};

  \node[prlab,above=0.4mm of a1] {$\pi{=}.9$};
  \node[prlab,above=0.4mm of a2] {$\pi{=}.9$};
  \node[prlab,above=0.4mm of a3] {$\pi{=}.9$};
  \node[prlab,below=0.4mm of a4] {$\pi{=}.1$};
  \node[prlab,below=0.4mm of a5] {$\pi{=}.1$};

  \draw[ent,relw=0.90] (a1) -- (a2) node[wlab,above]{$.90$};
  \draw[ent,relw=0.85] (a2) -- (a3) node[wlab,above]{$.85$};
  \draw[con,relw=0.90] (a3) -- (a4) node[wlab,left,pos=0.5]{$.90$};
  \draw[con,relw=0.85] (a3) -- (a5) node[wlab,right,pos=0.5]{$.85$};

  \node[font=\tiny,align=center,text=black!65,anchor=north] at (0.45,-0.90)
       {both conflicts point\\\emph{outward}, at the\\two false claims};
  % invisible strut matching panel B's annotation depth, so the two panels
  % share a baseline instead of one riding higher than the other.
  \path (0.45,-2.62) -- (0.45,0.98);
\end{tikzpicture}\\[3pt]
{\scriptsize \textbf{A:} a support chain over the true claims;
 each false claim \emph{attributed, then refused} from $a_3$.}
\end{minipage}
&
\begin{minipage}[b]{\linewidth}\centering
\begin{tikzpicture}[x=\mtvx, y=\mtvy]
  % Six edges over five nodes, so B needs more room than A: the two false
  % claims sit a full row lower and the row is wider, which keeps the
  % a4 -> a3 edge and the a2 -> a5 curve from crossing on top of each other.
  \node[atomtrue]  (a1) at (0,0)        {$a_1$};
  \node[atomtrue]  (a2) at (1.15,0)     {$a_2$};
  \node[atomtrue]  (a3) at (2.3,0)      {$a_3$};
  \node[atomfalse] (a4) at (0.55,-1.55) {$a_4$};
  \node[atomfalse] (a5) at (1.75,-1.55) {$a_5$};

  % Priors sit clear of the two orange change annotations.
  \node[prlab,above=0.4mm of a1] {$\pi{=}.9$};
  \node[prlab,above=0.4mm of a3] {$\pi{=}.9$};
  \node[prlab,above=0.4mm of a2] {$\pi{=}.9$};
  \node[prlab,left=0.5mm of a4]  {$\pi{=}.1$};
  \node[prlab,right=0.5mm of a5] {$\pi{=}.1$};

  \draw[ent,relw=0.90] (a1) -- (a2) node[wlab,above]{$.90$};
  \draw[con,relw=0.90] (a2) -- (a3) node[wlab,above]{$.90$};
  % the reversed edge: the false a4 now feeds the true a3
  \draw[ent,relw=0.88] (a4) -- (a3)
        node[wlab,above left=-0.4mm and -1.0mm,pos=0.30]{$.88$};
  \draw[con,relw=0.90] (a1) -- (a4) node[wlab,left,pos=0.55]{$.90$};
  \draw[ent,relw=0.85] (a4) -- (a5) node[wlab,below]{$.85$};
  % a2 -> a5 swings out to the right of a5 so it never lands on a4 -> a3
  \draw[con,relw=0.88] (a2) to[out=-8,in=25]
        node[wlab,right,pos=0.55,inner sep=1pt]{$.88$} (a5);

  % --- the two changes that carry the whole difference ---------------------
  % The two annotations are numbered and placed in the free space to the
  % RIGHT of the graph and BELOW it, so neither overlays a prior label, an
  % edge, or a weight. Each carries a dotted leader to what it marks.
  % (1) marks the a2--a3 edge, whose coupling flipped.
  \node[chgnum] (n1) at (2.98,-0.46) {1};
  \draw[chglead] (n1) -- (1.90,-0.06);
  % (2) rings a4, whose role reversed. The ring is drawn BEHIND the node so it
  % cannot cover the label; the number sits outside it.
  \begin{scope}[on background layer]
    \node[chg,circle,minimum size=9.4mm,inner sep=0pt] at (0.55,-1.55) {};
  \end{scope}
  \node[chgnum] (n2) at (-0.42,-1.05) {2};
  \draw[chglead] (n2) -- (0.22,-1.35);
  \node[chglab,anchor=north west,align=left] at (2.55,-0.75)
       {\textbf{1} support $\to$ \con{}:\\B denies its own\\conclusion $a_3$};
  \node[chglab,anchor=north,align=center] at (0.9,-2.20)
       {\textbf{2} the arrow \emph{reverses}: the false $a_4$\\is now a premise
        \emph{for} the true $a_3$};
\end{tikzpicture}\\[3pt]
{\scriptsize \textbf{B:} the same five claims, \emph{rearranged}. Half its
 conflicts point \emph{inward}, at what it asserts.}
\end{minipage}
\end{tabularx}

\medskip
%% ===== band (c): the per-claim posteriors ===============================
\begin{tikzpicture}[x=1mm, y=1mm]
  \def\bw{30}          % bar-track width
  \def\rh{5.2}         % row height
  \def\ca{47}          % left edge, track A
  \def\cb{88}          % left edge, track B

  \node[font=\bfseries\scriptsize,anchor=west] at (0,4.2)  {Claim};
  \node[font=\bfseries\scriptsize] at (\ca+0.5*\bw,4.2)    {$\Pr(a_i{=}1)$ under A};
  \node[font=\bfseries\scriptsize] at (\cb+0.5*\bw,4.2)    {$\Pr(a_i{=}1)$ under B};
  \draw[black!35,line width=0.3pt] (0,2.0) -- (127,2.0);

  % {i}{prior}{postA}{postB}{fill}{edge}{tag}{gloss} -- the colour is carried in
  % the list rather than branched on, so no \ifthenelse (not loaded) is needed.
  \foreach \i/\pri/\pa/\pb/\bc/\ec/\tag/\gloss in {%
      1/0.9/0.9386/0.9601/blue!35!white/blue!55!black/true/heliopause crossing\, Aug 2012,
      2/0.9/0.9891/0.9732/blue!35!white/blue!55!black/true/plasma-density rise after it,
      3/0.9/0.9920/0.5126/blue!35!white/blue!55!black/true/{\emph{therefore} in interstellar space},
      4/0.1/0.0129/0.0055/red!30!white/red!60!black/false/instrument failed pre-crossing,
      5/0.1/0.0199/0.0177/red!30!white/red!60!black/false/left the Solar System entirely}
  {
    \pgfmathsetmacro{\y}{-\i*\rh}
    \node[anchor=west,font=\scriptsize] at (0,\y)
         {$a_\i$ \textcolor{black!55}{(\tag)}\ \gloss};
    \foreach \x in {\ca,\cb}{
      \draw[black!22,line width=0.25pt] (\x,\y-1.7) rectangle (\x+\bw,\y+1.7);
    }
    \pgfmathsetmacro{\wa}{\pa*\bw}
    \pgfmathsetmacro{\wb}{\pb*\bw}
    \fill[\bc,draw=\ec,line width=0.3pt] (\ca,\y-1.7) rectangle (\ca+\wa,\y+1.7);
    \fill[\bc,draw=\ec,line width=0.3pt] (\cb,\y-1.7) rectangle (\cb+\wb,\y+1.7);
    % each claim's own prior, as a tick across the track
    \foreach \x in {\ca,\cb}{
      \pgfmathsetmacro{\tk}{\x+\pri*\bw}
      \draw[black!75,line width=0.6pt,dash pattern=on 0.6mm off 0.5mm]
            (\tk,\y-2.5) -- (\tk,\y+2.5);
    }
    \node[anchor=west,font=\tiny] at (\ca+\bw+1.4,\y) {\pa};
    \node[anchor=west,font=\tiny] at (\cb+\bw+1.4,\y) {\pb};
  }

  % the diagnostic callout: a3 under B
  \pgfmathsetmacro{\ythree}{-3*\rh}
  \draw[orange!85!black,line width=0.65pt,rounded corners=0.7mm]
        (\cb-0.9,\ythree-2.9) rectangle (\cb+\bw+9.5,\ythree+2.9);
  \draw[orange!85!black,line width=0.5pt,-{Stealth[length=1.3mm]}]
        (\cb+\bw+9.5,\ythree) -- (\cb+\bw+13,\ythree);

  \foreach \x in {\ca,\cb}{
    \node[font=\tiny,text=black!55] at (\x,-5*\rh-3.8)     {$0$};
    \node[font=\tiny,text=black!55] at (\x+\bw,-5*\rh-3.8) {$1$};
  }
  \node[font=\tiny,text=black!60,anchor=west] at (0,-5*\rh-3.8)
       {dashed rule $=$ that claim's own prior $\pi_i$};
\end{tikzpicture}

\smallskip
{\footnotesize\raggedright
\textcolor{orange!55!black}{\textbf{The diagnostic.}} $a_3$ is \textbf{true of
the world}, yet B drags it to a coin flip --- $0.992 \to 0.513$, far below its
own $0.9$ prior --- because its only support now arrives from $a_4$, a claim
the response itself denies. No aggregate score reveals this; a per-claim
marginal locates it.\par}

\medskip
%% ===== band (d): the scores =============================================
\begin{tabularx}{\textwidth}{@{}p{31mm}@{\hspace{8pt}}X@{}}
\toprule
\small\textbf{Coherence score}
& \small $\LCS_{\mm}$:\quad
  \textcolor{blue!55!black}{\textbf{$0.590$}}~(A)
  \ \ $\longrightarrow$\ \
  \textcolor{red!60!black}{\textbf{$0.494$}}~(B)
  \qquad $\boldsymbol{\Delta = -0.097}$ \\
\scriptsize\textcolor{black!60}{exact, by $2^5{=}32$-world enum.}
& \scriptsize $\LCS_{\cons}$: $0.963 \to 0.414$ \ \textbullet\
  $\log Z$: $-0.95 \to -3.78$ \ \textbullet\
  prior-only baseline $0.580$ (A \emph{above} it, B \emph{below}) \\
\midrule
\multicolumn{2}{@{}X@{}}{\scriptsize The priors enter \emph{only} through the
  unary factors, which are identical in the two responses. The entire
  $0.590\to0.494$ drop is produced by the pairwise relation factors --- that
  is, by \textbf{arrangement alone}.} \\
\bottomrule
\end{tabularx}

% scripts/make_motivating_png.sh defines \mtvnocaption to render the figure
% for standalone export, where the surrounding prose supplies the explanation.
% Undefined in the paper, so the caption below is typeset normally there.
\ifdefined\mtvnocaption\else
\caption{\textbf{Two responses, the same five claims, different coherence.}
  Five claims about Voyager~1 are fixed once and for all: $a_1,a_2,a_3$ true of
  the world (priors $\pi_i{=}0.9$, blue) and $a_4,a_5$ false ($\pi_i{=}0.1$,
  red). Both responses answer the \emph{same} query and assert \emph{all five}
  claims, so every per-claim factuality check returns the same verdict on both
  and no decompose-then-verify score can separate them. \textbf{A} attributes the two falsehoods and then refuses each
  on a stated ground; \textbf{B} asserts them and reasons \emph{from} them,
  offering the false $a_4$ as the reason the true $a_3$ holds while denying
  $a_3$ outright --- then denying $a_4$ and $a_5$ in its closing sentence.
  Middle: the relation graphs those spans induce (solid blue \ent{}, dashed red
  \con{}, line width growing with the mined weight $p$), drawn at one shared
  scale; the two numbered orange callouts \textbf{1} and \textbf{2} mark the
  only two changes, which carry the whole difference. Bottom: the posterior marginals of the coherence MRF, and the
  score. Mentioning a falsehood is not what costs B its score; asserting it and
  reasoning from it is. \S\ref{sec-example} works this pair formally and
  App.~\ref{app-examples} gives both responses in full.}
\label{fig:motivating}
\fi
\end{figure}
